\chapter{定義による拡大}\label{ch:definitions}

テキストに現れる「$\sim$ を $\cdots$ で定める」という定義は, 形式的には言語に新しい記号を追加し, その記号の意味を定める公理(定義公理)を理論に追加することである. 本章ではこの操作を定式化し, それが理論を本質的に強くしないこと(保存性)と, 新記号を含む論理式がもとの言語の論理式に翻訳できること(消去可能性)を示す. 本章は Kunen \cite{kunen2011} の定義 I.2.2, 定理 I.2.3, I.2.4 および Shoenfield \cite{shoenfield} \S 4.6 に従う.

以下 $T$ を言語 $\mathcal{L}$ の理論とする.

\section{述語記号と関数記号の定義}

\begin{dfn}[述語記号の定義]\label{dfn:pred-def}
  $\varphi(x_1, \ldots, x_n)$ を $\mathcal{L}$ の論理式で自由変数が $x_1,\ldots,x_n$ の中にあるものとし, $P$ を $\mathcal{L}$ にない $n$ 項関係記号とする. $\mathcal{L}' = \mathcal{L} \cup \{P\}$, $T' = T \cup \{\delta_P\}$ とおく. ここで
  \[ \delta_P :\equiv \forall x_1 \cdots \forall x_n\,(P x_1 \cdots x_n \leftrightarrow \varphi(x_1, \ldots, x_n)) \]
  である. $T'$ を $T$ の\textbf{述語記号 $P$ の定義による拡大}といい, $\delta_P$ を $P$ の\textbf{定義公理}という. $n = 0$ も許す(命題記号の定義).
\end{dfn}

\begin{dfn}[関数記号の定義]\label{dfn:func-def}
  $\varphi(x_1,\ldots,x_n,y)$ を $\mathcal{L}$ の論理式で
  \[ T \vdash \forall x_1 \cdots \forall x_n \exists ! y\, \varphi(x_1, \ldots, x_n, y) \tag{$\ast$}\]
  をみたすものとし, $F$ を $\mathcal{L}$ にない $n$ 項関数記号とする. $\mathcal{L}' = \mathcal{L} \cup \{F\}$, $T' = T \cup \{\delta_F\}$,
  \[ \delta_F :\equiv \forall x_1 \cdots \forall x_n\, \varphi(x_1, \ldots, x_n, F x_1 \cdots x_n) \]
  とおく. $T'$ を $T$ の\textbf{関数記号 $F$ の定義による拡大}といい, $\delta_F$ を $F$ の\textbf{定義公理}という. $n = 0$ のとき $F$ は定数記号である.
\end{dfn}

$(\ast)$ を確かめることが, 「$F x_1 \cdots x_n$ を $\varphi$ をみたす唯一の $y$ と定義する」ための\textbf{well-definedness の証明}である. 以下の章では, 関数記号の定義の前に必ず $(\ast)$ を定理として証明する.

\begin{lem}\label{lem:func-def-iff}
  関数記号 $F$ の定義による拡大 $T'$ において $T' \vdash \forall \vec{x}\forall y\,(y = F\vec{x} \leftrightarrow \varphi(\vec{x}, y))$.
\end{lem}
\begin{proof}
  ($\to$) $y = F\vec{x}$ と $\delta_F$ を $\forall$E したもの $\varphi(\vec{x}, F\vec{x})$ から Eq で $\varphi(\vec x, y)$. ($\leftarrow$) $(\ast)$ の一意性部分を $\forall$E して $\varphi(\vec x, y) \land \varphi(\vec x, F\vec x) \to y = F\vec x$ を得るので, $\varphi(\vec x, F\vec x)$ と TC. DT と Gen でまとめる.
\end{proof}

\section{保存性と消去可能性}

$T$ から定義による拡大を有限回(あるいは可算回)繰り返して得られる理論を, 単に $T$ の定義による拡大という. 各段階の定義公理の論理式は, その時点までに導入された記号を用いてよい.

\begin{dfn}[翻訳]\label{dfn:translation}
  $T'$ を $T$ の, 1つの記号の定義による拡大とする. $\mathcal{L}'$ の論理式 $\psi$ に対して, $\mathcal{L}$ の論理式 $\psi^\ast$ を次で定める.
  \begin{itemize}
    \item 新記号が述語記号 $P$ のとき: $\psi$ に現れる各原子論理式 $P t_1 \cdots t_n$ を $\varphi(t_1, \ldots, t_n)$ に置き換える(必要なら $\varphi$ の束縛変数を付け替える).
    \item 新記号が関数記号 $F$ のとき: $\psi$ が原子論理式で $F$ を含むならば, $\psi$ に現れる $F u_1 \cdots u_n$ の形の部分項であって $u_i$ が $F$ を含まないもの(最も内側の出現)を1つ選び, その出現をすべて新しい変数 $y$ で置き換えた原子論理式を $\psi_0$ として
    \[ \psi^\ast :\equiv \exists y\,(\varphi(u_1, \ldots, u_n, y) \land \psi_0)^\ast \]
    とする(右辺の $^\ast$ は $F$ の出現が減った論理式への再帰). $F$ を含まない原子論理式はそのままとし, $\neg, \land, \exists$ については構成に関して再帰的に定める.
  \end{itemize}
  何回かの拡大の場合は, 新しい記号から順に翻訳を合成する.
\end{dfn}

\begin{thm}[定義による拡大の基本定理 {\cite[定理 I.2.3]{kunen2011}}]\label{thm:defext}
  $T'$ を $T$ の定義による拡大とする.
  \begin{enumerate}[label=(\arabic*)]
    \item (消去可能性) $\mathcal{L}'$ の任意の論理式 $\psi$ について $T' \vdash \psi \leftrightarrow \psi^\ast$.
    \item (保存性) $\mathcal{L}$ の任意の論理式 $\psi$ について, $T' \vdash \psi$ ならば $T \vdash \psi$.
    \item 従って $\mathcal{L}'$ の任意の論理式 $\psi$ について, $T' \vdash \psi$ ならば $T \vdash \psi^\ast$.
  \end{enumerate}
\end{thm}
\begin{proof}
  1つの記号の拡大について示せば, 帰納法で一般の場合が従う.

  (1) $\psi$ の構成に関する帰納法. $\neg, \to$ の場合は帰納法の仮定から TC で, $\forall$ の場合は帰納法の仮定に Gen を施して $\forall x(\chi \leftrightarrow \chi^\ast)$ とし, (A3) と TC で $\forall x\chi \leftrightarrow \forall x \chi^\ast$ を得る. 原子論理式 $P\vec t$ の場合は $\delta_P$ の $\forall$E. 原子論理式 $\psi$ が $F$ を含む場合, 定義\ref{dfn:translation}の記法で $\psi \equiv (\psi_0)^y_{F\vec u}$ であり, 補題\ref{lem:func-def-iff}と等号公理 (A6) による書き換え(第\ref{sec:generator}節)から
  \[ T' \vdash \psi \leftrightarrow \exists y(\varphi(\vec u, y) \land \psi_0) \]
  が従う. 実際 ($\to$) は $\varphi(\vec u, F\vec u)$ と $\psi$ から $\exists$I, ($\leftarrow$) は $\varphi(\vec u, y)$ から $y = F\vec u$ を得て Eq. 右辺は $\psi$ より $F$ の出現の少ない論理式から作られるので, $F$ の出現数に関する帰納法で $(\psi_0)$ 部分に帰納法の仮定が使える.

  (2) 健全性定理と完全性定理 (Enderton 	\cite{enderton} \S 2.5) による. $\mathcal{A}$ を $T$ の任意のモデルとする. 述語記号の場合 $P^{\mathcal{A}} = \{\vec a \in A^n : \mathcal{A} \vDash \varphi[\vec a]\}$, 関数記号の場合 $F^{\mathcal{A}}(\vec a) = $ 「$\mathcal{A} \vDash \varphi[\vec a, b]$ となる唯一の $b$」(健全性定理より $\mathcal{A} \vDash (\ast)$ なので well-defined)と定めて $\mathcal{L}'$ 構造 $\mathcal{A}'$ を作ると, $\mathcal{A}' \vDash T'$ である. $T' \vdash \psi$ なら健全性定理より $\mathcal{A}' \vDash \psi$ であり, $\psi$ は $\mathcal{L}$ の論理式なので $\mathcal{A} \vDash \psi$ (充足関係は $\psi$ に現れない記号の解釈によらない). $\mathcal{A}$ は任意だから $T \vDash \psi$, 完全性定理より $T \vdash \psi$.

  (3) (1) より $T' \vdash \psi^\ast$ であり, $\psi^\ast$ は $\mathcal{L}$ の論理式なので (2) から従う.
\end{proof}

\begin{rem}
  (2) には, $T'$ からの証明を $T$ からの証明に変換する構文論的(有限的)な証明もある(Shoenfield \cite{shoenfield} \S 4.6, Kunen \cite{kunenFoM} \S II.15). 構文論的証明では, 証明木の各論理式を翻訳 $\psi \mapsto \psi^\ast$ し, 翻訳によって壊れた箇所(定義公理と, $F$ を含む項への (A2), (A6) の適用)を $(\ast)$ の証明で補う. 本ノートでは証明木は拡大された理論 $T'$ の中で書き, 定理\ref{thm:defext}によって $\ZFC$ の定理であることを保証する.
\end{rem}

\section{部分的にしか定義されない関数}

テキストの $\sup X$ や $\lim a_n$ のように, 引数がある条件をみたすときにだけ値が定まる関数がある. これを全域的な関数記号として導入するために, 次の約束をする. この約束は空集合 $\emptyset$ (定義\ref{dfn:emptyset})の導入後に用いる.

\begin{dfn}[記述による関数記号]\label{dfn:description}
  $T \vdash \forall \vec{x}\,\forall y\,\forall y'\,(\varphi(\vec x, y) \land \varphi(\vec x, y') \to y = y')$ (一意性のみ)が成り立つとき,
  \[ \hat\varphi(\vec x, y) :\equiv (\exists z\,\varphi(\vec x, z) \land \varphi(\vec x, y)) \lor (\neg\exists z\,\varphi(\vec x, z) \land y = \emptyset) \]
  とおくと $T \vdash \forall \vec x\exists ! y\,\hat\varphi(\vec x, y)$ であるから(場合分けによる. 後者の場合は $\emptyset$ が証拠), $\hat\varphi$ によって関数記号 $F$ を定義できる. このとき「$F\vec x$ を $\varphi(\vec x, y)$ をみたす唯一の $y$ と定める(そのような $y$ がなければ $\emptyset$)」という. 定義公理から直ちに
  \[ T' \vdash \forall \vec x\,(\exists y\,\varphi(\vec x, y) \to \varphi(\vec x, F\vec x)) \]
  が従う. 「存在しないとき $\emptyset$」という値は便宜的なもので, 以後の定理で使われることはない.
\end{dfn}

\section{内包記法}

集合 $a$ と論理式 $\varphi(x, \vec z)$ について, 分出公理と外延性公理から(第\ref{ch:zfc}章)
\[ \ZFC \vdash \forall a \forall \vec z\, \exists ! y\,\forall x\,(x \in y \leftrightarrow x \in a \land \varphi(x, \vec z)) \]
が成り立つ(定理\ref{thm:separation-term}). そこで論理式 $\varphi$ ごとに関数記号を導入し, その値を $\{x \in a \mid \varphi(x, \vec z)\}$ と書く. これは $a$ と $\vec z$ を引数とする関数記号であり, 定義公理は
\[ \forall a\forall \vec z\,\forall x\,(x \in \{x \in a \mid \varphi\} \leftrightarrow x \in a \land \varphi) \]
である. 一方, 集合であるとは限らない「クラス」$\{x \mid \varphi(x)\}$ は記号としては導入せず, $t \in \{x \mid \varphi\}$ を $\varphi(t)$ の略記としてのみ用いる.
